% DVote election lifecycle state machine.
% Bare tikzpicture: \input this inside a figure environment.
% The preamble supplies tikz with arrows.meta, positioning, shapes.geometric,
% fit, backgrounds and calc.
%
% Greyscale-safe conventions:
%   single outline = a phase the election passes through
%   double outline = a terminal phase reached by counting a result
%   dashed outline = the terminal phase reached by abandoning the election
%
% Each transition carries its call, its time gate and who may make it. The
% authority holds only configuration and cancellation, because a transition
% only it can make is also a transition it can withhold.
\begin{tikzpicture}[
    font=\footnotesize,
    phase/.style    = {draw=black, thick, rounded corners=3pt, fill=white,
                       align=center, inner sep=3pt,
                       minimum width=4.6cm, minimum height=1.0cm},
    terminal/.style = {phase, double, double distance=1.1pt},
    aborted/.style  = {phase, dashed, fill=black!7,
                       minimum width=3.0cm},
    trans/.style     = {-{Latex[length=2.6mm,width=2.2mm]}, draw=black,
                       line width=1.1pt},
    abort/.style    = {-{Latex[length=2.2mm,width=1.8mm]}, draw=black!65,
                       semithick, dash pattern=on 2.5pt off 2pt,
                       rounded corners=3pt},
    tlab/.style     = {align=left, inner sep=2pt, text width=6.9cm,
                       font=\footnotesize,
                       execute at begin node={\hyphenpenalty=10000\exhyphenpenalty=10000}},
  ]

  % ==== the phases ====
  \node[phase]    (created)  at (6.0,0)
    {\texttt{Created}\\cloned by the factory,\\candidates and windows set};
  \node[phase]    (regn)     at (6.0,-2.8)
    {\texttt{Registration}\\identity commitments collected};
  \node[phase]    (voting)   at (6.0,-5.6)
    {\texttt{Voting}\\root frozen, ballots accepted};
  \node[phase]    (tallying) at (6.0,-8.4)
    {\texttt{Tallying}\\ballots refused; the counts\\already stand in storage};
  \node[terminal] (final)    at (6.0,-11.2)
    {\texttt{Finalized}\\result written, immutable};
  \node[aborted]  (cancelled) at (1.5,-1.4)
    {\texttt{Cancelled}\\terminal, no result};

  % ==== the forward path ====
  \draw[trans] (created.south)  -- (regn.north);
  \draw[trans] (regn.south)     -- (voting.north);
  \draw[trans] (voting.south)   -- (tallying.north);
  \draw[trans] (tallying.south) -- (final.north);

  \node[tlab, anchor=west] at (8.6,-1.4)
    {\texttt{openRegistration()}, time-gated: \texttt{block.timestamp} must have
     reached the announced opening, after which anyone may call it. A transition
     only the authority can make is one the authority can also withhold.};
  \node[tlab, anchor=west] at (8.6,-4.2)
    {\texttt{openVoting()}, time-gated on the voting window and likewise callable
     by anyone. It also requires a non-empty leaf-set CID beside the root, so a
     registrar that withholds the published roll stops the election rather than
     going unnoticed. The root was written when registration closed and cannot
     change from here.};
  \node[tlab, anchor=west] at (8.6,-7.0)
    {\texttt{closeVoting()}, time-gated on the voting window closing, and
     callable by anyone once that timestamp has passed, so the authority
     cannot stop the count.};
  \node[tlab, anchor=west] at (8.6,-9.8)
    {\texttt{finalize()}, likewise callable by anyone, the authority
     included. Writes the result and emits \texttt{Finalized}; no later call
     can change it.};

  % ==== the one abandonment path ====
  \draw[abort] (created.west) -- (3.4,0) -- (3.4,-0.9) -- (cancelled.north east);
  \draw[abort] (regn.west)    -- (3.4,-2.8) -- (3.4,-1.9) -- (cancelled.south east);
  \node[align=left, anchor=north, text width=3.3cm, inner sep=2pt,
        execute at begin node={\hyphenpenalty=10000\exhyphenpenalty=10000}] at (1.5,-2.2)
    {\texttt{cancel()}, the authority only, and only before voting opens. After
     that a ballot may exist, and cancelling the election would be a veto over
     a result the authority can already see.};

\end{tikzpicture}
